% Universidade Estadual de Feira de Santana
% Modelo de Proposta de Dissertação de Mestrado em Ciencia da Computacao
% Adaptado por: José Amancio e Roberto Bittencourt, UEFS
% Trabalho original feito por: Ana Cristina Alves de Oliveira (cristina@dsc.ufcg.edu.br)
% Orientador: Francisco Brasileiro
% em Fevereiro de 2020
% Agosto de 2019

% E necessario o arquivo "algorithmic.sty", se for construir algoritmos

\documentclass[a4paper,titlepage,12pt]{article}

\usepackage{color}
\usepackage{longtable}
\usepackage[brazil,english]{babel}
\usepackage[utf8]{inputenc}
% \usepackage[latin1]{inputenc}

\usepackage{times}
\usepackage[T1]{fontenc}
\usepackage{fancyheadings}
\usepackage{fancyvrb}

\usepackage{algorithmic}
\usepackage[nothing]{algorithm}
\usepackage{latexsym}

\usepackage{graphicx,url}

% para uso de referência direta e indireta, com \citet, \citep, etc.
\usepackage[round]{natbib}

\usepackage[pdfborder={0 0 0},breaklinks=true,bookmarksopen=true,bookmarksopenlevel=1]{hyperref}
\usepackage{pdflscape}
\usepackage{lipsum}

\sloppy

% Comandos de estilo e espacamento ----------------------------------------
\newlength{\defbaselineskip}
\setlength{\defbaselineskip}{\baselineskip}
\newcommand{\setlinespacing}[1]%
           {\setlength{\baselineskip}{#1 \defbaselineskip}}

\setcounter{topnumber}{2}
\renewcommand{\topfraction}{.7}
\setcounter{bottomnumber}{1}
\renewcommand{\bottomfraction}{.3}
\setcounter{totalnumber}{3}
\renewcommand{\textfraction}{.2}
\renewcommand{\floatpagefraction}{.5}
\setcounter{dbltopnumber}{2}
\renewcommand{\dbltopfraction}{.7}
\renewcommand{\dblfloatpagefraction}{.5}
%
\oddsidemargin -28pt
\evensidemargin -28pt
\marginparwidth 50pt
\marginparsep 5pt
\topmargin -27pt
\hoffset 15mm
\textheight 237mm
\textwidth 155mm
\renewcommand{\baselinestretch}{1.5}
%

% ------------------------------------------------------------------------

%%%%%  Ambiente de redefinicoes para construcao de Algoritmos
% Apenas usar se precisar fazer algoritmos

\renewcommand{\algorithmicend}{\textbf{fim}}
\renewcommand{\algorithmicwhen}{\textbf{quando}}
\renewcommand{\algorithmicprimitive}{\textbf{in?cio}}
\renewcommand{\algorithmicendprimitive}{\textbf{fim}}
\renewcommand{\algorithmicendwhen}{\textbf{fim}}
\renewcommand{\algorithmicif}{\textbf{se}}
\renewcommand{\algorithmicthen}{\textbf{ent\~{a}o}}
\renewcommand{\algorithmicelse}{\textbf{sen\~{a}o}}
\renewcommand{\algorithmicendif}{\textbf{fim-se}}
\renewcommand{\algorithmicelsif}{\algorithmicelse\ \algorithmicif}
\renewcommand{\algorithmicendif}{\algorithmicend\ \algorithmicif}
\renewcommand{\algorithmicfor}{\textbf{para}}
\renewcommand{\algorithmicforall}{\textbf{paratodos}}
\renewcommand{\algorithmicdo}{\textbf{fa\c{c}a}}
\renewcommand{\algorithmicendfor}{\algorithmicend\ \algorithmicfor}
\renewcommand{\algorithmicwhile}{\textbf{enquanto}}
\renewcommand{\algorithmicendwhile}{\algorithmicend\ \algorithmicwhile}
\renewcommand{\algorithmicloop}{\textbf{loop}}
\renewcommand{\algorithmicendloop}{\algorithmicend\ \algorithmicloop}
\renewcommand{\algorithmicrepeat}{\textbf{repita}}
\renewcommand{\algorithmicuntil}{\textbf{at\'{e}}}
\renewcommand{\algorithmicwaituntil}{\textbf{espera at?}}
\floatname{algorithm}{Algoritmo}

%%%%%%

% Color definitions (RGB model)
\definecolor{mycolor1}{rgb}{0.753,0.753,0.753}

\begin{document}

% Primeira Folha do Documento %%%%%%%%%%%%%%%%%%%%%%%%%%%%%%%%%%%%%%%%%%%%%

\pagestyle{empty}

\begin{center}
{\textbf{\Large \textsc{Universidade Estadual de Feira de Santana}}}
\end{center}

\begin{center}
{\Large \textsc{\textbf{Programa de Pós-Graduação em \\Ciência da Computação}}}
\end{center}

~\\ \\

\begin{center}
{\LARGE \textsc{\textbf{Proposta de Dissertação de Mestrado}}}
\end{center}

~\\ \\

\begin{center}
{\Large \textsc{\textbf{<<Título da Proposta>>}}}
\end{center}

~\\ \\

\begin{center}
\textbf{\textsc{Mestrando(a)} \\
\textsc{<<Seu Nome>>}}
\end{center}

\begin{center}
\textbf{\textsc{Orientador(a)} \\
\textsc{<<Nome(s) do(s) Orientador(es)>>}}
\end{center}

~\\ \\ \\

\begin{center}
\textbf{{\large \textsc{Feira de Santana}}
\\
{\large \textsc{<<Mes>> - <<Ano>>}}}
\end{center}

\newpage
\cleardoublepage

%%%%%%%%%%%%%%%%%%%%%%%%%%%%%%%%%%%%%%%%%%%%%%%%%%%%%%%%%%%%%%%%%%%%%%%%%%%%%%%%

\selectlanguage{brazil}
% \selectlanguage{english}

% ConFigura os números das páginas para algarismos romanos
\pagestyle{plain}
\pagenumbering{roman}

%%%%%%%%%%%%%%%%%%%%%%%%%%%%%%%%%%%%%%%%%%%%%%%%%%%%%%%%%%%%%%%%%%%%%%%%%%%%%%%%

\section*{Resumo}
\addcontentsline{toc}{section}{Resumo}

\vspace{3cm}

Escreva aqui o texto do resumo de sua proposta em português. \lipsum[1]

\clearpage

%%%%%%%%%%%%%%%%%%%%%%%%%%%%%%%%%%%%%%%%%%%%%%%%%%%%%%%%%%%%%%%%%%%%%%%%%%%%%%%%


\tableofcontents
\newpage


% Corpo do documento -------------------

% Configura os números das páginas para algarismos indo-arabicos
\pagestyle{plain}
\setcounter{page}{1}
\pagenumbering{arabic}

%%%%%%%%%%%%%%%%%%%%%%%%%%%%%%%%%%%%%%%%%%%%%%%%%%%%%%%%%%%%%%%%%%%%%%%%%%%%%%%%
%% Definicao do cabecalho: secao do lado esquerdo e numero da pagina do lado direito
\pagestyle{fancy}
\addtolength{\headwidth}{\marginparsep}\addtolength{\headwidth}{\marginparwidth}\headwidth
= \textwidth
%\renewcommand{\sectionmark}[1]{\markboth{#1}{}}
\renewcommand{\sectionmark}[1]{\markright{\thesection\ #1}}\lhead[\fancyplain{}{\bfseries\thepage}]%
         {\fancyplain{}{\emph{\rightmark}}}\rhead[\fancyplain{}{\bfseries\leftmark}]%
             {\fancyplain{}{\bfseries\thepage}}\cfoot{}

%%%%%%%%%%%%%%%%%%%%%%%%%%%%%%%%%%%%%%%%%%%%%%%%%%%%%%%%%%%%%%%%%%%%%%%%%%%%%%%%

\input{source/a_introducao.tex}
\input{source/b_fundamentos.tex}
\input{source/c_metodologia.tex}
\input{source/d_cronograma.tex}
\input{source/e_resultados.tex}

%%%%%%%%%%%%%%%%%%%%%%%%%%%%%%%%%%%%%%%%%%%%%%%%%%%%%%%%%%%%%%%%%%%%%%%%%%%%%%%%
%% BIbliografia
\newpage

\bibliographystyle{apalike-portuguese} % note o uso do pacote natbib (para uso de referencia direta e indireta com \citet, \citep, etc.). 
% \bibliographystyle{apalike} % Caso a proposta seja em inglês, descomentar para usar o estilo apalike original que vem com o pacote natbib. 
\bibliography{referencias} % arquivos com as entradas bib.

%%%%%%%%%%%%%%%%%%%%%%%%%%%%%%%%%%%%%%%%%%%%%%%%%%%%%%%%%%%%%%%%%%%%%%%%%%%%%%%%

\newpage

\input{source/z_apendices.tex}

\end{document}
